% sec8_2.tex — Section 8.2: Truncated Regression Models

A \textbf{limited dependent variable (LDV) model} is a regression model in which either
the dependent variable $y_t$ is constrained in some way or the observations for which
$y_t$ does not meet some prespecified criterion are excluded from the sample. The
LDV model of the former sort is called a \textbf{censored regression model}, while the
latter is called a \textbf{truncated regression model}. This section presents truncated
regression models; censored regression models are presented in the next section.

\subsection*{The Model}

Suppose that $\{y_t, \mathbf{x}_t\}$ is i.i.d.\ satisfying
\begin{equation}
  y_t = \mathbf{x}_t'\boldsymbol{\beta}_0 + \varepsilon_t, \quad
  \varepsilon_t \mid \mathbf{x}_t \sim N(0, \sigma_0^2).
  \label{eq:8.2.1} \tag{8.2.1}
\end{equation}
The model would be just the standard linear regression model with normal errors,
were it not for the following feature: only those observations that meet some prespecified rule for $y_t$ are included in the sample. We will consider only the simplest
truncation rule: $y_t > c$ where $c$ is a known constant. This rule is sometimes called
a ``truncation from below.'' Once this case is worked out, it should not be hard to
consider more general truncation rules.

\subsection*{Truncated Distributions}

To proceed, we need two results from probability theory.

\medskip
\noindent\textit{Density of a Truncated Random Variable:}\quad
If a continuous random variable $y$ has a density function $f(y)$ and $c$ is a constant, then its distribution after
the truncation $y > c$ is defined over the interval $(c, \infty)$ and is given by
\begin{equation}
  f(y \mid y > c) = \frac{f(y)}{\text{Prob}(y > c)}.
  \label{eq:8.2.2} \tag{8.2.2}
\end{equation}

\begin{figure}[ht]
  \centering
  % Placeholder for Figure 8.1: Effect of Truncation
  % Shows density of N(0,1) (solid) and density after truncation (dotted)
  % with truncation point c marked on horizontal axis
  \includegraphics[width=0.9\textwidth,trim=50 350 50 350,clip]{figures/fig_8_1.png}
  \caption{Effect of Truncation}
  \label{fig:8.1}
\end{figure}

Figure~\ref{fig:8.1} illustrates how truncation alters the distribution for $N(0,1)$. The solid
curve is the density of $N(0,1)$. Because it is a density, the area under the curve is
unity. The dotted curve is the distribution after truncation, which is defined over
the interval $(c, \infty)$. It lies above the solid curve over this interval so that the area
under the dotted curve is unity.

The second result follows.

\medskip
\noindent\textit{Moments of the Truncated Normal Distribution:}\quad
If $y \sim N(\mu_0, \sigma_0^2)$ and $c$ is a constant, then the mean and the variance of the truncated distribution are
\begin{align}
  \E(y \mid y > c) &= \mu_0 + \sigma_0 \lambda(v),
  \label{eq:8.2.3} \tag{8.2.3} \\
  \Var(y \mid y > c) &= \sigma_0^2 \{1 - \lambda(v)[\lambda(v) - v]\},
  \label{eq:8.2.4} \tag{8.2.4}
\end{align}
where $v = (c - \mu_0)/\sigma_0$ and
$\lambda(v) \equiv \frac{\phi(v)}{1 - \Phi(v)}$.

This formula for the mean clearly shows that the sample mean of draws from the
truncated distribution is not consistent for $\mu_0$. This is an example of the \textbf{sample
selection bias}. The bias equals $\sigma_0 \lambda(v)$. The function $\lambda(v)$ is called the
\textbf{inverse Mill's ratio} or the \textbf{hazard function}. The function $\lambda(v)$ is convex and asymptotes
to $v$ as $v \to \infty$ and to zero as $v \to -\infty$. So its derivative $\lambda'(v)$ is between 0 and
1. It is easy to show that
\begin{equation}
  \lambda'(v) = \lambda(v)(\lambda(v) - v).
  \label{eq:8.2.5} \tag{8.2.5}
\end{equation}

Using \eqref{eq:8.2.3} and \eqref{eq:8.2.4}, we can show how truncation alters the form of
regression of $y_t$ on $\mathbf{x}_t$. By \eqref{eq:8.2.1}, the distribution of $y_t$ conditional on $\mathbf{x}_t$ \emph{before}
truncation is $N(\mathbf{x}_t'\boldsymbol{\beta}_0, \sigma_0^2)$. Observation~$t$ is in the sample if and only if $y_t > c$. So
\begin{align}
  \E(y_t \mid \mathbf{x}_t,\, t \text{ in the sample})
  &= \mathbf{x}_t'\boldsymbol{\beta}_0
     + \sigma_0 \lambda\!\Bigl(\frac{c - \mathbf{x}_t'\boldsymbol{\beta}_0}{\sigma_0}\Bigr),
  \label{eq:8.2.6} \tag{8.2.6} \\
  \Var(y_t \mid \mathbf{x}_t,\, t \text{ in the sample})
  &= \sigma_0^2 \biggl\{1 - \lambda\!\Bigl(\frac{c - \mathbf{x}_t'\boldsymbol{\beta}_0}{\sigma_0}\Bigr)
     \biggl[\lambda\!\Bigl(\frac{c - \mathbf{x}_t'\boldsymbol{\beta}_0}{\sigma_0}\Bigr)
     - \frac{c - \mathbf{x}_t'\boldsymbol{\beta}_0}{\sigma_0}\biggr]\biggr\}.
  \label{eq:8.2.7} \tag{8.2.7}
\end{align}
This formula shows that the OLS estimate of the $\mathbf{x}_t$ coefficient from the linear
regression of $y_t$ on $\mathbf{x}_t$ is not consistent because the correction term
$\sigma_0 \lambda\bigl(\frac{c - \mathbf{x}_t'\boldsymbol{\beta}_0}{\sigma_0}\bigr)$,
which is correlated with $\mathbf{x}_t$, would be included in the error term in the linear regression. Since the functional form of the hazard function $\lambda(v)$ is known (under the
normality assumption for $\varepsilon_t$), we could avoid the sample selection bias by applying nonlinear least squares to estimate $(\boldsymbol{\beta}_0, \sigma_0^2)$, but maximum likelihood should
be the preferred estimation method because it is asymptotically more efficient.

In passing, we note that the sample selection bias does not arise if selection is
based on the regressors, not on the dependent variable. Suppose that observation~$t$
is in the sample if $\phi(\mathbf{x}_t) \geq 0$. Then
\begin{align}
  \E(y_t \mid \mathbf{x}_t,\, t \text{ in the sample})
  &= \E[\mathbf{x}_t'\boldsymbol{\beta}_0 + \varepsilon_t
     \mid \mathbf{x}_t,\, \phi(\mathbf{x}_t) \geq 0] \notag \\
  &= \mathbf{x}_t'\boldsymbol{\beta}_0
     + \E[\varepsilon_t \mid \mathbf{x}_t,\, \phi(\mathbf{x}_t) \geq 0] \notag \\
  &= \mathbf{x}_t'\boldsymbol{\beta}_0.
  \label{eq:8.2.8} \tag{8.2.8}
\end{align}
The last equality holds since
$\E[\varepsilon_t \mid \mathbf{x}_t,\, \phi(\mathbf{x}_t) \geq 0]
= \E(\varepsilon_t \mid \mathbf{x}_t) = 0$ if $\phi(\mathbf{x}_t) \geq 0$.

\subsection*{The Likelihood Function}

Derivation of the likelihood function utilizes \eqref{eq:8.2.2}. As just noted, the pretruncation distribution of $y_t \mid \mathbf{x}_t$ is
$N(\mathbf{x}_t'\boldsymbol{\beta}_0, \sigma_0^2)$, whose density is
\begin{equation}
  \frac{1}{\sqrt{2\pi\sigma_0^2}}
  \exp\biggl[-\frac{1}{2}\Bigl(\frac{y_t - \mathbf{x}_t'\boldsymbol{\beta}_0}{\sigma_0}\Bigr)^2\biggr]
  = \frac{1}{\sigma_0}\,\phi\!\Bigl(\frac{y_t - \mathbf{x}_t'\boldsymbol{\beta}_0}{\sigma_0}\Bigr),
  \label{eq:8.2.9} \tag{8.2.9}
\end{equation}
where $\phi(\cdot)$ is the density of $N(0,1)$. The probability that observation~$t$ is in the
sample is given by
\begin{align}
  \text{Prob}(y_t > c \mid \mathbf{x}_t)
  &= 1 - \text{Prob}(y_t \leq c \mid \mathbf{x}_t) \notag \\
  &= 1 - \text{Prob}\!\biggl(\frac{y_t - \mathbf{x}_t'\boldsymbol{\beta}_0}{\sigma_0}
     \leq \frac{c - \mathbf{x}_t'\boldsymbol{\beta}_0}{\sigma_0}
     \;\bigg|\; \mathbf{x}_t\biggr) \notag \\
  &= 1 - \Phi\!\Bigl(\frac{c - \mathbf{x}_t'\boldsymbol{\beta}_0}{\sigma_0}\Bigr)
     \quad\bigl(\text{since } \tfrac{y_t - \mathbf{x}_t'\boldsymbol{\beta}_0}{\sigma_0}
     \mid \mathbf{x}_t \sim N(0,1)\bigr).
  \label{eq:8.2.10} \tag{8.2.10}
\end{align}
Therefore, by \eqref{eq:8.2.2}, the post-truncation density, defined over $(c, \infty)$, is
\begin{equation}
  f(y_t \mid \mathbf{x}_t,\, t \text{ in the sample};\,
    \boldsymbol{\beta}_0, \sigma_0^2)
  = \frac{\frac{1}{\sigma_0}\,\phi\!\bigl(\frac{y_t - \mathbf{x}_t'\boldsymbol{\beta}_0}{\sigma_0}\bigr)}
    {1 - \Phi\!\bigl(\frac{c - \mathbf{x}_t'\boldsymbol{\beta}_0}{\sigma_0}\bigr)}.
  \label{eq:8.2.11} \tag{8.2.11}
\end{equation}
This is the density (conditional on $\mathbf{x}_t$) of observation~$t$ \emph{in the sample}. Taking logs
and replacing $(\boldsymbol{\beta}_0, \sigma_0^2)$ by their hypothetical values $(\boldsymbol{\beta}, \sigma^2)$, we obtain the log conditional likelihood for observation~$t$:
\begin{multline}
  \log f(y_t \mid \mathbf{x}_t; \boldsymbol{\beta}, \sigma^2) \\
  = \biggl\{-\frac{1}{2}\log(2\pi) - \frac{1}{2}\log(\sigma^2)
    - \frac{1}{2}\Bigl(\frac{y_t - \mathbf{x}_t'\boldsymbol{\beta}}{\sigma}\Bigr)^2\biggr\}
    - \log\biggl[1 - \Phi\!\Bigl(\frac{c - \mathbf{x}_t'\boldsymbol{\beta}}{\sigma}\Bigr)\biggr].
  \label{eq:8.2.12} \tag{8.2.12}
\end{multline}
Were it not for the last term, this would be the log likelihood (conditional on $\mathbf{x}_t$) for
the usual linear regression model. The last term is needed because the value of the
dependent variable for observation~$t$, having passed the test $y_t > c$ and thus being
in the sample, is never less than or equal to~$c$.

\subsection*{Reparameterizing the Likelihood Function}

This log likelihood function can be made easier to analyze if we introduce the
following reparameterization:
\begin{equation}
  \boldsymbol{\delta} = \boldsymbol{\beta}/\sigma, \quad \gamma = 1/\sigma.
  \label{eq:8.2.13} \tag{8.2.13}
\end{equation}
(We have considered this reparameterization in Example~7.7 for the linear regression model.) This is a one-to-one mapping between $(\boldsymbol{\beta}, \sigma^2)$ and
$(\boldsymbol{\delta}, \gamma)$, with the inverse mapping given by
$\boldsymbol{\beta} = \boldsymbol{\delta}/\gamma$, $\sigma^2 = 1/\gamma^2$.
The reparameterized log conditional likelihood is
\begin{multline}
  \log \tilde{f}(y_t \mid \mathbf{x}_t; \boldsymbol{\lambda}) \\
  = \bigl[-\tfrac{1}{2}\log(2\pi) + \log(\gamma)
    - \tfrac{1}{2}(\gamma y_t - \mathbf{x}_t'\boldsymbol{\delta})^2\bigr]
    - \log[1 - \Phi(\gamma c - \mathbf{x}_t'\boldsymbol{\delta})].
  \label{eq:8.2.14} \tag{8.2.14}
\end{multline}

The objective function in ML estimation is $1/n$ times the log likelihood of the
sample. For a random sample, the log likelihood of the sample is the sum over~$t$
of the log likelihood for observation~$t$. Thus, the objective function in the ML
estimation of the reparameterized truncated regression model, denoted
$\tilde{Q}_n(\boldsymbol{\delta}, \gamma)$, is the average over~$t$ of \eqref{eq:8.2.14}. The ML estimator
$(\hat{\boldsymbol{\delta}}, \hat{\gamma})$ of $(\boldsymbol{\delta}_0, \gamma_0)$ is the $(\boldsymbol{\delta}, \gamma)$
that maximizes this objective function.

\subsection*{Verifying Consistency and Asymptotic Normality}

The next task is to derive the score and the Hessian and check the conditions
for consistency and asymptotic normality for the reparameterized log likelihood
\eqref{eq:8.2.14}. (On the first reading, you may wish to skip the discussion on consistency
and asymptotic normality.)

\subsubsection*{Score and Hessian for Observation \texorpdfstring{$t$}{t}}

A tedious calculation yields
\begin{equation}
  \underset{((K+1) \times 1)}{\mathbf{s}(\mathbf{w}_t;\, \boldsymbol{\delta}, \gamma)}
  = \begin{bmatrix}
      (\gamma y_t - \mathbf{x}_t'\boldsymbol{\delta})\mathbf{x}_t \\[4pt]
      \frac{1}{\gamma} - (\gamma y_t - \mathbf{x}_t'\boldsymbol{\delta})y_t
    \end{bmatrix}
  + \lambda(v_t)
    \begin{bmatrix} -\mathbf{x}_t \\[4pt] c \end{bmatrix},
  \label{eq:8.2.15} \tag{8.2.15}
\end{equation}
\begin{multline}
  \underset{((K+1) \times (K+1))}{\mathbf{H}(\mathbf{w}_t;\, \boldsymbol{\delta}, \gamma)}
  = -\begin{bmatrix}
      \mathbf{x}_t\mathbf{x}_t' & -y_t\mathbf{x}_t \\[4pt]
      -y_t\mathbf{x}_t' & \frac{1}{\gamma^2} + y_t^2
    \end{bmatrix} \\
  + \lambda(v_t)[\lambda(v_t) - v_t]
    \begin{bmatrix}
      \mathbf{x}_t\mathbf{x}_t' & -c\,\mathbf{x}_t \\[4pt]
      -c\,\mathbf{x}_t' & c^2
    \end{bmatrix},
  \label{eq:8.2.16} \tag{8.2.16}
\end{multline}
where $K$ is the number of regressors,
$\mathbf{w}_t = (y_t, \mathbf{x}_t')'$, and
$v_t \equiv \gamma c - \mathbf{x}_t'\boldsymbol{\delta}$ $\bigl(= \frac{c - \mathbf{x}_t'\boldsymbol{\beta}}{\sigma}\bigr)$.
The Hessian is \emph{not} everywhere negative semidefinite even after the reparameterization. So the objective function $\tilde{Q}_n(\boldsymbol{\delta}, \gamma)$ is not concave.\footnote{However, it can be shown that the Hessian of the objective function $Q_n(\boldsymbol{\theta})$ is negative semidefinite at any
solution to the likelihood equations (see Orme, 1989). Since this rules out any saddle points or local minima,
having more than two local maxima is impossible. Therefore, provided that the maximum of
$\tilde{Q}_n(\boldsymbol{\delta}, \gamma)$ occurs at
a point interior to the parameter space, the first-order condition is necessary and sufficient for a \emph{global} maximum.
If an algorithm such as Newton--Raphson produces convergence, the point of convergence is the global maximum.}
The relevant consistency theorem, therefore, is Proposition~7.5, not Proposition~7.6, while the relevant
asymptotic normality theorem remains Proposition~7.9.

\subsubsection*{Consistency}

Condition~(a) of Proposition~7.5 is the (conditional) density identification
\begin{align*}
  &\tilde{f}(y_t \mid \mathbf{x}_t;\, \boldsymbol{\delta}, \gamma)
   \neq \tilde{f}(y_t \mid \mathbf{x}_t;\, \boldsymbol{\delta}_0, \gamma_0) \\
  &\qquad\text{with positive probability for }
   (\boldsymbol{\delta}, \gamma) \neq (\boldsymbol{\delta}_0, \gamma_0),
\end{align*}
where $(\boldsymbol{\delta}_0, \gamma_0)$ is the true parameter value and the parameterized log density
$\tilde{f}(y_t \mid \mathbf{x}_t;\, \boldsymbol{\delta}, \gamma)$ is given in \eqref{eq:8.2.14}. Clearly, given
$(y_t, \mathbf{x}_t)$, the value of
$\tilde{f}(y_t \mid \mathbf{x}_t;\, \boldsymbol{\delta}, \gamma)$
is different for different values of $\gamma$. So identification boils down to the condition
that $\mathbf{x}_t'\boldsymbol{\delta} \neq \mathbf{x}_t'\boldsymbol{\delta}_0$ with positive probability for
$\boldsymbol{\delta} \neq \boldsymbol{\delta}_0$. But \eqref{eq:8.1.10} indicates that
this condition is satisfied if $\E(\mathbf{x}_t\mathbf{x}_t')$ is nonsingular. Turning to condition~(b) of the
proposition (the dominance condition for
$\tilde{f}(y_t \mid \mathbf{x}_t;\, \boldsymbol{\delta}, \gamma)$), it is easy to show that
the condition is satisfied under the nonsingularity of
$\E(\mathbf{x}_t\mathbf{x}_t')$.\footnote{Only interested readers: what needs to be shown is
$\E[\sup_{\boldsymbol{\theta} \in \Theta} |\log \tilde{f}(y_t \mid \mathbf{x}_t; \boldsymbol{\theta})|] < \infty$,
where $\boldsymbol{\theta} = (\boldsymbol{\delta}', \gamma)'$ and $\Theta$ is a compact set.
An extension of the argument in Example~7.8 shows that there is a dominating
function for the bracketed term in \eqref{eq:8.2.14}. The inequality (7.2.14) can be used to show that there is a dominating
function for $\log[1 - \Phi(cy - \mathbf{x}_t'\boldsymbol{\delta})]$.}
Therefore, the ML
estimator $(\hat{\boldsymbol{\delta}}, \hat{\gamma})$ of $(\boldsymbol{\delta}_0, \gamma_0)$ is consistent under the nonsingularity of $\E(\mathbf{x}_t\mathbf{x}_t')$.

\subsubsection*{Asymptotic Normality}

A rather tedious calculation using \eqref{eq:8.2.6}, \eqref{eq:8.2.7}, and \eqref{eq:8.2.13} shows that
$\E[\mathbf{s} \mid \mathbf{x}_t] = 0$. An extremely tedious calculation yields the conditional information
matrix equality: $\E[\mathbf{s}\mathbf{s}' \mid \mathbf{x}_t] = -\E[\mathbf{H} \mid \mathbf{x}_t]$.
Therefore, condition~(3) of Proposition~7.9 is satisfied. We will not discuss verification of conditions~(4) and~(5) of
the proposition.\footnote{There seems to be no published work that verifies conditions~(4) and~(5). For the case where $\{\mathbf{x}_t\}$ is a
sequence of fixed constants, Sapra (1992) proved asymptotic normality under the assumption that $\mathbf{x}_t$ is bounded
and $\lim_{n\to\infty} \frac{1}{n}\sum_{t=1}^{n} \mathbf{x}_t\mathbf{x}_t'$ is nonsingular. (His proof is for the more general setting of serially correlated observations.) It would seem reasonable to conjecture that for the case where $\mathbf{x}_t$ is random (as here), a sufficient
condition for consistency and asymptotic normality is the nonsingularity of $\E(\mathbf{x}_t\mathbf{x}_t')$.}

Thus, we conclude: if $\{y_t, \mathbf{x}_t\}$ is i.i.d.\ and $\E(\mathbf{x}_t\mathbf{x}_t')$ is nonsingular, then the ML
estimator $(\hat{\boldsymbol{\delta}}, \hat{\gamma})$ of the truncated regression model is consistent and asymptotically
normal with the asymptotic variance consistently estimated by
\begin{equation}
  \widehat{\avar}(\hat{\boldsymbol{\delta}}, \hat{\gamma})
  = -\biggl[\frac{1}{n}\sum_{t=1}^{n}
    \mathbf{H}(\mathbf{w}_t;\, \hat{\boldsymbol{\delta}}, \hat{\gamma})\biggr]^{-1},
  \label{eq:8.2.17} \tag{8.2.17}
\end{equation}
where $\mathbf{w}_t = (y_t, \mathbf{x}_t')'$ and $\mathbf{H}(\mathbf{w}_t; \boldsymbol{\theta})$ is given in \eqref{eq:8.2.16}.

\subsection*{Recovering Original Parameters}

By the invariance property of ML (see Section~7.1), the ML estimate of
$(\boldsymbol{\beta}_0, \sigma_0^2)$
is $\hat{\boldsymbol{\beta}} = \hat{\boldsymbol{\delta}}/\hat{\gamma}$ and
$\hat{\sigma}^2 = 1/\hat{\gamma}^2$, the value of the inverse mapping. Clearly, if the ML
estimate of the reparameterized model, $(\hat{\boldsymbol{\delta}}, \hat{\gamma})$, is consistent, then so is the ML estimate $(\hat{\boldsymbol{\beta}}, \hat{\sigma}^2)$ thus recovered. The asymptotic variance of
$(\hat{\boldsymbol{\beta}}, \hat{\sigma}^2)$ can be obtained
by the delta method (see Lemma~2.5) as
$\hat{\mathbf{J}}\,\widehat{\avar}(\hat{\boldsymbol{\delta}}, \hat{\gamma})\,\hat{\mathbf{J}}'$,
where $\hat{\mathbf{J}}$ is the estimate of the Jacobian matrix for the inverse mapping
($\boldsymbol{\beta} = \boldsymbol{\delta}/\gamma$, $\sigma^2 = 1/\gamma^2$) evaluated at
$(\hat{\boldsymbol{\delta}}, \hat{\gamma})$:
\begin{equation}
  \hat{\mathbf{J}} = \begin{bmatrix}
    \frac{1}{\hat{\gamma}}\mathbf{I} & -\frac{\hat{\boldsymbol{\delta}}}{\hat{\gamma}^2} \\[6pt]
    \mathbf{0}' & -\frac{2}{\hat{\gamma}^3}
  \end{bmatrix}.
  \label{eq:8.2.18} \tag{8.2.18}
\end{equation}

\bigskip
\noindent\rule{\textwidth}{0.4pt}
\textbf{QUESTIONS FOR REVIEW}
\medskip

\begin{enumerate}
\item (Second moment of $y_t$)\quad
  From \eqref{eq:8.2.6} and \eqref{eq:8.2.7}, derive:
  \[
    \E(y_t^2 \mid \mathbf{x}_t,\, t \text{ in sample})
    = \frac{1}{\gamma^2}\bigl[1 + \lambda(v_t)\gamma c
      + (\mathbf{x}_t'\boldsymbol{\delta})^2
      + \lambda(v_t)\mathbf{x}_t'\boldsymbol{\delta}\bigr],
  \]
  where $v_t \equiv \gamma c - \mathbf{x}_t'\boldsymbol{\delta}$.

\item (Estimation of truncated regression model by GMM)\quad
  The truncated regression model can also be estimated by GMM. Consider the following $K + 1$
  orthogonality conditions. The first $K$ of them are
  \begin{equation}
    \underset{(K \times 1)}{\mathbf{g}_1(\mathbf{w}_t;\, \boldsymbol{\delta}, \gamma)}
    = \biggl(y_t - \frac{1}{\gamma}\mathbf{x}_t'\boldsymbol{\delta}
      - \frac{\lambda(v_t)}{\gamma}\biggr) \cdot \mathbf{x}_t,
    \label{eq:8.2.19} \tag{8.2.19}
  \end{equation}
  where $\mathbf{w}_t = (y_t, \mathbf{x}_t')'$ and
  $v_t \equiv \gamma c - \mathbf{x}_t'\boldsymbol{\delta}$. The $(K+1)$-th orthogonality condition is
  \begin{equation}
    g_2(\mathbf{w}_t;\, \boldsymbol{\delta}, \gamma)
    = y_t^2 - \frac{1}{\gamma^2}\bigl[1 + \lambda(v_t)\gamma c
      + (\mathbf{x}_t'\boldsymbol{\delta})^2
      + \lambda(v_t)\mathbf{x}_t'\boldsymbol{\delta}\bigr].
    \label{eq:8.2.20} \tag{8.2.20}
  \end{equation}

  \begin{enumerate}[label=(\alph*)]
  \item Verify that
    $\E[\mathbf{g}_1(\mathbf{w}_t; \boldsymbol{\delta}_0, \gamma_0)] = \mathbf{0}$ and
    $\E[g_2(\mathbf{w}_t; \boldsymbol{\delta}_0, \gamma_0)] = 0$.
    \textbf{Hint:} Use the Law of Total Expectations. $v_t$ is a function of $\mathbf{x}_t$.

  \item Show that $(\boldsymbol{\delta}^*, \gamma^*)$ satisfies the likelihood equations if and only if it satisfies
    the $K+1$ orthogonality conditions. \textbf{Hint:} Let $s_2(\mathbf{w}_t; \boldsymbol{\delta}, \gamma)$ be the $(K+1)$-th
  \end{enumerate}
\end{enumerate}
